\documentclass[11pt]{article}
\usepackage[margin=1in]{geometry}
\usepackage{amsmath,amssymb}
\usepackage{booktabs}
\usepackage{hyperref}

\title{Scenario 1 --- Multi-Day Receding-Horizon MPC\\
       for Mobile Charging Station Dispatch\\[2pt]
       \large Formal Model of the Cross-Day Receding Horizon (Approach 2 tree)}
\author{Sections 1--11: the deterministic MILP (Approach 0's per-day baseline, and the
per-scenario building block of Approach 2's stochastic extension). Section~\ref{sec:stochastic}:
the stochastic (scenario-based) extension itself --- Approach 2's actual closed loop.}
\date{}

\newcommand{\dt}{\Delta t}

\begin{document}
\maketitle

Formal companion to the code in \texttt{Receding\_Horizon/code/}: the cross-day
window MILP is built in \texttt{3\_MCSModel.jl} (\texttt{build\_window\_model}); the
overnight recharge is scheduled inside that same MILP, not a separate phase; the multi-day
closed loop is in \texttt{4\_MPCLoop.jl} (\texttt{run\_mpc}). Plain-English descriptions
and the data schema are in \texttt{README.md}; a line-by-line code audit is in
\texttt{constraints\_code\_vs\_model.txt}. Sections~\ref{sec:sets}--\ref{sec:plant} below
document the \textbf{deterministic} model: solved every 15 minutes by Approach~0's per-day
baseline replay (\emph{once per day}, then replayed open-loop for that day, to measure the
value of re-planning) and reused --- unchanged, per scenario --- as the building block of
Approach~2's own stochastic model in Section~\ref{sec:stochastic}, which is what this tree's
actual multi-day closed loop (\texttt{run\_mpc}) solves every 15 minutes. The daytime window
MILP itself is identical either way, up to the scenario-linking described in
Section~\ref{sec:stochastic}.

The controller is a \textbf{multi-day, cross-day RECEDING horizon}. The loop simulates
$D_{\text{tot}} = n_{\text{days}}+1$ days ($n_{\text{days}}$ reported plus one dropped
\emph{buffer} day). Each day is a \textbf{daytime block} of $n_{\text{day}}$ intervals
($n_{\text{day}}$ inferred from the last work-available interval $+$ a return buffer,
capped at the full day $n_{\text{int}}=96$). Every 15-min step re-solves a MILP over a
\textbf{fixed-length, exactly-24h rolling window}
\[
  K_{\text{win}} = \bigl[\,k_0,\ \dots,\ k_0 + n_{\text{day}} - 1\,\bigr],
\]
always $n_{\text{day}}$ intervals ahead of \emph{now}, sliding forward by one interval each
step; only the first interval is applied and the
window recedes. \textbf{Nights ARE intervals}: the day-block is the full 24\,h
($n_{\text{day}}=n_{\text{int}}=96$), so the overnight MCS recharge is scheduled by the MILP
itself, driven by the time-of-use price. There is no reset bridge and no separate deterministic
recharge phase. Battery state and any unfinished work carry over continuously. The buffer day
absorbs the energy-neutral wrap-up so it never lands on a reported day.

The activity powers are a \textbf{fixed} calibrated estimate (Fork B); the MILP plans on
the mean $\mu$, and the stochastic plant draws realized powers from a pool pre-sampled from
$\mathcal{N}(\mu,\sigma)$, one draw per (excavator, activity) occurrence (shared across every
approach that simulates the plant, so their realizations are comparable), truncated at $0$
(idle $\sigma=0$). The model is not re-fit online.

\paragraph{Window geometry.} For a global interval index $k$, let
$\mathrm{wd}(k)=((k-1)\bmod n_{\text{day}})+1$ be its position \emph{within} its day and
$\mathrm{dayof}(k)=\lfloor(k-1)/n_{\text{day}}\rfloor+1$ its day. Daily profiles (price,
carbon, work availability) repeat each day and are indexed by $\mathrm{wd}(k)$. The
day-boundary index $k_{\text{term}}=d_0\,n_{\text{day}}$ (for a window whose first day is $d_0$)
is the last interval before the next day-start; $b_{\text{term}}=k_{\text{term}}+1$ is the boundary
where both terminal targets are pinned. The window may see a short tail past $k_{\text{term}}$.

%======================================================================
\section{Sets and indices}
\label{sec:sets}
\begin{tabular}{ll}
$m \in \mathcal{M}$ & mobile charging stations (MCS) \\
$e \in \mathcal{E}$ & construction EVs (excavators) \\
$i,j \in \mathcal{N}$ & nodes; $\mathcal{N}_g$ grid, $\mathcal{N}_c$ sites \\
$a \in \mathcal{B}$ & activities $=\{\text{dig},\text{load},\text{trv},\text{idle}\}$ \\
$k$ & global interval in the cross-day window $K_{\text{win}}$ \\
$\mathrm{wd}(k),\mathrm{dayof}(k)$ & within-day index / day of $k$ \\
$k_{\text{term}},b_{\text{term}}$ & last interval before / boundary at the next day-start \\
$\mathcal{D}$ & day-blocks present in the window \\
$\mathcal{T}_b$ & interval boundaries of the window \\
\end{tabular}

\medskip
$n_{\text{day}}=n_{\text{int}}=96$ is the \textbf{full 24\,h day-block} at $\dt=0.25$\,h;
the overnight hours are ordinary intervals of the same MILP. $A_{i,e}=1$
iff CEV $e$ works at site $i$. Work availability $R_{i,e,k}=R^{\text{day}}_{i,e,\mathrm{wd}(k)}$
is $0$ outside the shift, but the MCS/CEVs may charge and the MCS may travel at any
daytime interval.

%======================================================================
\section{Parameters (selected)}
\begin{tabular}{ll}
$\dt$ & interval length (h), $0.25$ \\
$\lambda^{\text{buy}}_{k},\lambda^{\text{CO2}}_{k}$ & price (\$/kWh) / carbon intensity \\
$\overline{SOE}_\bullet,\underline{SOE}_\bullet,SOE^{\text{ini}}_\bullet$ & energy cap / floor / start (MCS, CEV) \\
$CH_m,DCH_m,DCH^{\text{plug}}_m,C^{\text{plug}}_m,\eta_m,CH^{\text{CEV}}_e$ & MCS / CEV charge-discharge parameters \\
$p_a$ & activity power (kW); fixed calibrated mean, MILP plans on it \\
$H^{\text{dig}}_{i,d},H^{\text{load}}_{i,d}$ & \textbf{per-day} dig / load quota (h) at site $i$, day $d$ \\
$r^{\text{dig}}_i,r^{\text{load}}_i$ & work still remaining at site $i$ carried into the window \\
$\tau_{i,j},k_{\text{trv}},R_{i,e,k}$ & travel time / energy / work availability \\
$\rho_{\text{miss}},\rho_{\text{lab}},\lambda^{\text{NC}},\lambda^{\text{OP}}$ & penalties \\
$\text{scale},\kappa=4,r$ & precedence ratio / travels-per-work ($\kappa_{\text{wt}}$) / rest cap \\
\end{tabular}

%======================================================================
\section{Decision variables}
\textbf{Continuous ($\ge0$):} $P^{\text{ch}}_{m,i,k}$, $P^{\text{dch}}_{m,i,k}$,
$P^{\text{MCS-CEV}}_{m,i,e,k}$, $P^{\text{work}}_{i,e,k}$, totals
$P^{\text{ch,tot}}_{m,k},P^{\text{dch,tot}}_{m,k}$, travel energy $L_{m,i,j,k}$, state
$SOE^{\text{MCS}}_{m,t},SOE^{\text{CEV}}_{e,t}$, peaks $P^{\text{NC}},P^{\text{OP}}$, and slacks
$s^{\text{miss}}_{i,a}$ (per site \& activity --- \emph{not} per day-block). This is the only slack
in the model: every other constraint is hard, and there are no \texttt{soft\_*} levers.
The objective sums $s^{\text{miss}}$ over all four activities although only dig and load appear in
a balance; the travel/idle entries are free $\ge0$ variables carrying positive cost, so they are
driven to $0$ (the source PDF's objective has the same quirk).

\smallskip
\textbf{Binary:} $u_{e,i,a,k}$, $\mu_{i,e,k}$, $\rho_{m,i,e,k}$, $z_{m,i,k}$,
$g^{\text{ch}}_{m,i,k}$, $x_{m,i,j,k}$, $y_{m,i,j,k}$, $\beta^{\text{arr}},\beta^{\text{dep}}$.

%======================================================================
\section{Objective (Eq.\ 1)}
\begin{align}
\min\ &\sum_{m,k}\lambda^{\text{buy}}_{\mathrm{wd}(k)}P^{\text{ch,tot}}_{m,k}\dt
      +\sum_{m,k}\tfrac{c_{\text{ton}}}{1000}\lambda^{\text{CO2}}_{\mathrm{wd}(k)}P^{\text{ch,tot}}_{m,k}\dt
      +\rho_{\text{miss}}\!\!\sum_{i\in\mathcal{N}_c,d\in\mathcal{D}}\!\!\bigl(s^{\text{miss,dig}}_{i,d}+s^{\text{miss,load}}_{i,d}\bigr)\notag\\
     &+\lambda^{\text{NC}}P^{\text{NC}}+\lambda^{\text{OP}}P^{\text{OP}}
      +\rho_{\text{lab}}\dt\!\!\sum_{m,i,j,k}\!\! y_{m,i,j,k}
      + W_{\text{p}}\!\sum_{i,k}\! s^{\text{prec}}_{i,k}
      + W_{\text{c}}\!\sum_{e,k}\!(s^{\text{pace}+}_{e,k}{+}s^{\text{pace}-}_{e,k})
      + W_{\text{t}}\!\sum_{e,d}\! s^{\text{term}}_{e,d}
      +10^{-6}\!\!\sum_{i\in\mathcal{N}_c,\,e,\,idx}\!\! idx\cdot \mu_{i,e,K_{idx}}.
\end{align}
The last term is CHANGE 2's tie-break penalty, added ONLY to this deterministic objective, NOT to
the stochastic window model's objective -- Approach 2's stochastic hedge already addresses
margin-for-error via scenarios (Change 4).

\paragraph{Change 2 --- earlier-charging tie-break (Issue 2).} Penalizes later CEV-charging (via
the CEV-acceptance indicator $\mu_{i,e,k}$, NOT $g^{\text{ch}}$ -- the diagnosed stall was the MCS
idling on discharge \emph{to the CEV} with plenty of its own charge to spare, gated by $\mu$;
$g^{\text{ch}}$, the MCS's own grid-charging indicator, was already 0 throughout that stall) at
weight $10^{-6}$, breaking ties toward earlier action without ever overriding a genuine cost
difference. $idx$ is the LOCAL position within the current re-solve window.

\paragraph{Change 3 --- terminal SOE\textsuperscript{CEV} shortfall penalty (Issue 1), NOT part of
this objective.} Evaluated AFTER the run from the realized trajectory at the kept-day boundary:
$\text{shortfall}_e=\max(SOE^{\text{CEV,ini}}_e-SOE^{\text{CEV,end}}_e,0)$, converted to hours by
the day's required weighted work-power rate (the full dig/load quota, not what was actually
achieved), charged as
$\rho_{\text{miss}}\times\text{shortfall\_hours}$ directly onto reported TOTAL cost.

\paragraph{Change 4 --- stratified scenario sampling (Issue 3).} \texttt{sample\_scenarios} now
draws 5 FIXED, evenly spread bins instead of 5 i.i.d.\ Normal draws, guaranteeing tail coverage
every re-solve. See the identical description in the Shrinking-Horizon sibling document.

\paragraph{Change 5 --- \texttt{n\_day\_run}.} The day-count knob (previously a hardcoded 2-day
schedule with two DIFFERENT quotas) is renamed \texttt{n\_day\_run}, defaulting to 1. For
\texttt{n\_day\_run}~$>1$, every day uses the SAME work requirement and starts from the previous
day's REAL ending state (day-to-day carry-over already existed here; only the default and
same-work semantics are new).

Price/carbon use the within-day index $\mathrm{wd}(k)$. The overnight recharge is scheduled
\emph{inside} this MILP, so its cost is already carried by the first two terms; there is no
separate overnight accounting.



%======================================================================
\section{Constraints}

\paragraph{Power aggregation \& flow (HARD).}
$P^{\text{ch,tot}}_{m,k}=\sum_{i\in\mathcal{N}_g}P^{\text{ch}}_{m,i,k}$,
$P^{\text{dch,tot}}_{m,k}=\sum_{i\in\mathcal{N}_c}P^{\text{dch}}_{m,i,k}$;
$P^{\text{dch}}_{m,i,k}=\sum_e P^{\text{MCS-CEV}}_{m,i,e,k}\le DCH_m z_{m,i,k}$;
$P^{\text{MCS-CEV}}_{m,i,e,k}\le DCH^{\text{plug}}_m\rho_{m,i,e,k}$;
$\sum_m P^{\text{MCS-CEV}}_{m,i,e,k}\le CH^{\text{CEV}}_e\mu_{i,e,k}$.
Grid exclusivity: $P^{\text{ch}}_{m,i,k}\le CH_m g^{\text{ch}}_{m,i,k}$,
$g^{\text{ch}}\le z$, $\sum_m g^{\text{ch}}_{m,i,k}\le 1$.

\paragraph{State of energy (HARD).}
\begin{align}
SOE^{\text{MCS}}_{m,k+1}&=SOE^{\text{MCS}}_{m,k}+\eta_m P^{\text{ch,tot}}_{m,k}\dt-\tfrac{1}{\eta_m}P^{\text{dch,tot}}_{m,k}\dt-L^{\text{tot}}_{m,k},\\
SOE^{\text{CEV}}_{e,k+1}&=SOE^{\text{CEV}}_{e,k}+\textstyle\sum_{m,i}P^{\text{MCS-CEV}}_{m,i,e,k}\dt-\sum_i P^{\text{work}}_{i,e,k}\dt,\\
\underline{SOE}_\bullet&\le SOE_\bullet\le\overline{SOE}_\bullet .
\end{align}
Both recursions run \emph{unbroken} across the whole window, nights included. There is no
overnight reset bridge and no separate deterministic recharge phase: because the day-block is the
full 24\,h ($n_{\text{day}}=n_{\text{int}}=96$), the overnight recharge is scheduled by the MILP
itself, driven by the time-of-use price in the objective.

\paragraph{Terminal energy targets (Eq.\ 8a / 8b, HARD).}
Both are pinned to $b_{\text{term}}=k_{\text{term}}+1$, the boundary at the \emph{next} day-start
($k_{\text{term}}=d_0\,n_{\text{day}}$ for a window whose first day is $d_0$). They are applied
unconditionally --- the window is one day-block long, so that boundary always lies inside it:
\begin{align}
SOE^{\text{MCS}}_{m,b_{\text{term}}}&= SOE^{\text{MCS,ini}}_m && \text{(8a, exact; recharge scheduled inside this MILP)},\\
SOE^{\text{CEV}}_{e,b_{\text{term}}}&\ge SOE^{\text{CEV,ini}}_e && \text{(8b, floor; overcharging allowed)}.
\end{align}
The CEV target is a \emph{floor} rather than Avik's exact equality: a CEV cannot discharge, so
under the stochastic plant an equality can become unrecoverable once a CEV drifts high; the floor
keeps the terminal reachable while guaranteeing the fleet ends at least as charged as it began.
Both constraints are \textbf{hard}, with no tolerance band and no soft variant.

\paragraph{Plugging, presence \& routing (Eq.\ 10, HARD).}
$\rho\le A$, $\rho\le z$, $\sum_e \rho_{m,i,e,k}\le C^{\text{plug}}_m$. Presence partition
$\sum_i z_{m,i,k}+\sum_{i\ne j}y_{m,i,j,k}=1$; $\beta^{\text{dep}}_{m,i,k}=\sum_j x_{m,i,j,k}$;
$\beta^{\text{arr}}$ from finishing trips (or a carried-in arrival);
$\beta^{\text{arr}}-\beta^{\text{dep}}=z_k-z_{k-1}$; $\beta^{\text{arr}}+\beta^{\text{dep}}\le1$;
flow balance (generalised for the carried-in start). Trip on $i\!\to\!j$ lasts $\tau_{i,j}$ and
costs $L=k_{\text{trv}}\dt\,y$. \textbf{(10e)} MCS parked at a grid node at the day-boundary:
$\sum_{i\in\mathcal{N}_g}z_{m,i,k_{\text{term}}}=1$.

\paragraph{Activity scheduling (Eq.\ 11, HARD).}
$\sum_a u_{e,i,a,k}=A_{i,e}$; $u\le A$; $\mu_{i,e,k}\le u_{e,i,\text{idle},k}$;
$P^{\text{work}}_{i,e,k}\le R_{i,e,k}A_{i,e}(1-\mu_{i,e,k})$;
$P^{\text{work}}_{i,e,k}=\sum_a p_a u_{e,i,a,k}$ (Eq.\ 5e), with $p_{\text{idle}}=0$.
(The code uses this 4-activity encoding; Avik's equivalent 3-activity form is
$\sum_a u\le1$, $\sum_a u+\mu\le1$.)

\paragraph{Work quota (Eq.\ 12c, SOFT shortfall + implicit HARD cap).}
Work is issued as a \textbf{per-day} schedule ($H^{\text{dig}}_{i,d}$, $H^{\text{load}}_{i,d}$),
but the roll-over is handled \emph{outside} the MILP: the closed loop adds each morning's quota to
a running backlog, $r^{\bullet}_i \mathrel{+}= H^{\bullet}_{i,d}$, and subtracts realized work as it
is applied. The MILP therefore sees only that backlog, and enforces one balance per site over the
intervals up to the day-boundary $k_{\text{term}}$:
\begin{equation}
\dt\!\!\sum_{e,\;k\in K_{\text{win}},\,k\le k_{\text{term}}}\!\!u_{e,i,\text{dig},k}
\;+\;s^{\text{miss}}_{i,\text{dig}}\;=\;\max(r^{\text{dig}}_i,0),
\end{equation}
and identically for loading. There is a \textbf{single} slack per (site, activity) --- it is
\emph{not} indexed by day-block --- and a \textbf{single} constraint per site, not one per day.

The shortfall $s^{\text{miss}}$ is soft (priced at $\rho_{\text{miss}}$ in Eq.~1), but since
$s^{\text{miss}}\ge0$ the equality also implies the \textbf{hard} upper bound
$\dt\sum u\le\max(r^{\bullet}_i,0)$: a CEV cannot do more work than the backlog still outstanding,
which is how ``no working ahead'' is obtained --- implicitly, and only up to $k_{\text{term}}$.
Intervals in the window's tail beyond $k_{\text{term}}$ are excluded from this balance entirely and
carry no quota accounting; they are unconstrained rather than budgeted (benign in practice, since
work costs energy and earns nothing until the following day's backlog is issued).

\paragraph{Precedence (Eq.\ 12d, HARD).}
Raw interval counts, \textbf{cumulative over the whole run} --- the counters do \emph{not} restart
each morning. The seed $C^{\text{load}}_i,C^{\text{dig}}_i$ is the realized work applied in every
earlier interval since the run began (the shared history is never reset), converted hours
$\to$ interval counts:
\begin{equation}
\frac{C^{\text{load}}_i}{\dt}+\!\!\sum_{e,\,\min K_{\text{win}}\le\tau\le k}\!\!u_{e,i,\text{load},\tau}
\ \le\ \text{scale}\left(\frac{C^{\text{dig}}_i}{\dt}+\!\!\sum_{e,\,\min K_{\text{win}}\le\tau\le k}\!\!u_{e,i,\text{dig},\tau}\right).
\end{equation}
Hard; there is no precedence slack. The seed aggregates \emph{per site}, whereas travel pacing
below seeds \emph{per CEV}; the two agree while $A_{i,e}$ is a one-to-one assignment.

\paragraph{Rest rule (Eq.\ 12e, HARD).}
With $r=\mathrm{round}(t_{\text{rest}}/\dt)$ (=4 at the defaults $t_{\text{rest}}=1$\,h,
$\dt=0.25$\,h; the code rounds, it does not take a ceiling), over any $(r{+}1)$-window inside
$K_{\text{win}}$,
$\sum_{k'=k}^{k+r}\sum_{a\in\{\text{dig,load,trv}\}}u_{e,i,a,k'}\le r$, \emph{plus a seam} seeded
from the applied Work/Break history so a work-run cannot leak across the every-15-min re-solves.
The counters do \emph{not} restart each morning; the history runs continuously. Because the
day-block is the full 24\,h, the overnight intervals sit in that history as idle, so a night
supplies its own break and the seam is harmless in practice.

\paragraph{Travel pacing (Eq.\ 13, HARD, no tolerance).}
With $\kappa=\kappa_{\text{wt}}=4$, for each assigned $(i,e)$, on cumulative travel $V$ vs
cumulative useful work $W$ (dig+load) --- both raw \emph{applied interval counts} off the $u$
indicator, not hours, seeded from \texttt{cum\_trv\_cnt\_e}/\texttt{cum\_work\_cnt\_e}:
\begin{equation}
W(k)-\kappa\ \le\ \kappa\,V(k)\ \le\ W(k).
\end{equation}
Hard; there are no pacing slacks and no tolerance is needed --- a battery-shortage-capped interval
(e.g.\ 14 of 15 minutes realized) still counts as one full travel/work interval, so the fractional
rounding residue that used to strand the band on an integer boundary (the old $\tau =
\texttt{pacing\_tol} = 0.05$ fix) cannot occur here.

Unlike precedence and the rest rule, the pacing seed is \textbf{scoped to the current calendar day
only}: \texttt{hist} is never reset across days, so the count is taken from
\texttt{today\_start} $= (\text{firstday}-1)\cdot n_{\text{day}} + 1$ onward --- only intervals
applied since this calendar day's boundary. A CEV's travel/work balance resets fresh every
morning. Applies to both the deterministic and scenario-linked stochastic builder.

\paragraph{Demand peaks (E1, HARD).}
$P^{\text{NC}}\ge$ carried-in peak and $\ge\sum_m P^{\text{ch,tot}}_{m,k}$ for all $k$;
$P^{\text{OP}}$ likewise on on-peak $k$ only.

%======================================================================
\section{Closed loop, overnight Phase 2 \& graceful fallback}
The MILP above is re-solved every 15 minutes over the receding cross-day window; only interval
$k_0$ is applied, then the real state (SOE, MCS position including mid-drive trips, per-site
remaining work, demand peaks, and the current day's applied-activity history) is advanced and the
window slides. The activity powers $p_a$ are the mean of a \emph{fixed} calibrated estimate
$\mathcal{N}(\mu,\sigma)$; the plant draws each realized power from a pool pre-sampled from
$\mathcal{N}(\mu,\sigma)$ (truncated at $0$, idle pinned to $0$) once, before the run, and
consumes the next unused sample for a $(\text{CEV},a)$ pair each time that pair actually occurs
--- not a fresh draw every interval; the model is not re-fit online.

\paragraph{No overnight phase.} The day-block is the full 24\,h, so the overnight recharge is an
ordinary part of the window MILP, scheduled against the time-of-use price and pinned by the
terminal equality $SOE^{\text{MCS}}_{m,b_{\text{term}}}=SOE^{\text{MCS,ini}}_m$. There is no
separate closed-form recharge routine.

\paragraph{No graceful fallback.} Every constraint except $s^{\text{miss}}$ is hard and there is no
soft re-solve. If a window is infeasible the plant simply \emph{holds} state for that interval and
the event is counted (\texttt{n\_infeasible}); nothing is relaxed and no second solve is attempted.

\paragraph{Buffer day.} The loop simulates $D_{\text{tot}}=n_{\text{days}}+1$ days and drops the
last; KPIs and figures use only days $1,\dots,n_{\text{days}}$, so the terminal wrap-up never lands
on a reported day.

%======================================================================
\section{Plant modes: what the disturbance actually is}
\label{sec:plant}
The MILP is identical in every mode; only what the \emph{plant} does with the applied plan
differs. Both \texttt{run\_mpc} (Approach~1) and \texttt{run\_one\_shot} (Approach~0) take a
\texttt{plant} switch:

\medskip
\begin{tabular}{ll}
\texttt{:sampled} & realized power $\hat p_{e,a}$ = the next unused draw from the pre-sampled\\
                  & pool; the within-interval activity split may be randomised. Realized\\
                  & $\ne$ planned, so the day drifts. \emph{(default)}\\[4pt]
\texttt{:mean}    & realized power $\hat p_{e,a}=\mu_a$, the same mean the MILP planned on, and\\
                  & each interval realizes its single planned activity in full. Realized\\
                  & $=$ planned \emph{exactly} within each day. No pool sample is consumed\\
                  & and no cursor advances, so a \texttt{:mean} run cannot perturb a\\
                  & \texttt{:sampled} run sharing the same pool, and needs no seed.\\
\end{tabular}

\medskip
Approach~0 under \texttt{:mean} is therefore the per-day MILPs' \emph{own optima}, chained through
the real end-of-day carry-over --- the deterministic reference. Exactly one mode runs per call
(\texttt{approach0\_plant}); Approach~1 is always \texttt{:sampled}. Which one is selected decides
what the reported gap means:
\begin{align*}
\text{A0}(\texttt{:sampled})\to\text{A1}\ &=\ \text{value of intra-day re-planning (like-for-like plant)},\\
\text{A0}(\texttt{:mean})\to\text{A1}\ &=\ \text{drift and re-planning together}.
\end{align*}

\paragraph{Scope of the disturbance.} Only the CEV side is stochastic. The MCS state is advanced
by the \emph{planned} value $SOE^{\text{MCS}}_{m,g_0+1}$ from the solved window, so grid draw,
travel loss and hence every cost term follow the plan exactly; the disturbance enters through CEV
depletion and reaches the objective only indirectly, through what the next re-solve decides to do
about it.

\paragraph{State guard.} The realized dig/load/travel duration for each CEV is capped by the
energy actually available before hitting $\underline{SOE}_e$: if the sampled draw would push the
balance below the floor, the executed duration is scaled down to exactly what the floor allows,
crediting only the affordable fraction of that interval's work (the remainder becomes idle) before
the SOE update is applied. This is a physical correction, not a numerical guard — no energy is
created or destroyed, and \texttt{rem\_dig}/\texttt{rem\_load} are only reduced by what was truly
completed. A residual \texttt{clamp} to $[\underline{SOE}_e,\overline{SOE}_e]$ remains afterward as
a pure safety net and should not bind under normal operation. Capping events are counted and
reported per run ($\texttt{n\_capped}$), because they are invisible to the MILP's own feasibility
status.

%======================================================================
\section{Approach 2: the stochastic (scenario-based) extension}
\label{sec:stochastic}
Everything above is the deterministic rolling-window MILP exactly as Approach~1 (and
Approach~0's per-day baseline) solve it, implemented unchanged by \texttt{build\_window\_model}
in \texttt{3\_MCSModel.jl}. Approach~2's controller solves a sibling model,
\texttt{build\_window\_model\_stochastic}, in its place at every re-solve. The formal
difference is \emph{identical} to the Shrinking-Horizon codebase's own \texttt{math\_model.tex},
\S13
(see that document for the full development and \texttt{Understanding\_Stochastic\_MPC.md} for
the conceptual explanation and a worked example) — only the window geometry differs, exactly
as it does between the two codebases' deterministic models above: the rolling window's own
day-block indexing $wd(k)$, $dayof(k)$ and the fixed next-boundary terminal targets
$k_{\text{term}}$, $b_{\text{term}}$ carry over into every scenario copy unchanged, in place of
the shrinking window's own terminal-at-day-end condition.

\paragraph{Scenario set and objective.} At each re-solve the controller samples $S$ power
vectors $\{\mu^{(s)}\}_{s=1}^S \sim \mathcal N(\mu,\sigma)^+$ (clipped at 0, equal weight
$\pi_s=1/S$, $S=5$ default) from a dedicated RNG stream disjoint from the plant's. Every
variable in \S1--\S6 gains a scenario index; every constraint holds per scenario, HARD; the
objective is $\min \sum_s \pi_s J^{(s)}$ with $\mu\to\mu^{(s)}$ substituted into the work-power
identity (Eq.~5e) of each scenario's copy — identical in form to the Shrinking-Horizon model.

\paragraph{Non-anticipativity.} Let $g_0=\text{first}(K)$ be the rolling window's own current
global interval (the ``here-and-now,'' called $k_0$ in the Shrinking-Horizon document). The
same control family
$\{u,\mu^{\text{plug}},\rho,z,g^{ch},x^{\text{trip}},y^{trv},\beta^{arr},\beta^{dep},
P^{ch}_{MCS},P^{MCS}_{CEV},P^{dch}_{MCS},P^{ch}_{tot},P^{dch}_{tot},L^{trv},L^{trv}_{tot}\}$
is tied, $X^{(s)}_{\dots,g_0}=X^{(1)}_{\dots,g_0}$ for $s=2,\dots,S$; $P^{work}_{i,e,g_0}$ and
$SOE^{CEV}_{e,g_0+1}$ are deliberately \emph{not} tied, for exactly the reason given in the
Shrinking-Horizon document: the chosen activity is shared, its energy cost is not. $s^{\text{miss}}$,
$P^{peak}_{NC}$, $P^{peak}_{OP}$ remain free window-wide aggregates. At $S=1$ this collapses
exactly to \S1--\S6's deterministic model with $\mu^{(1)}=\mu$.

\end{document}
